\chapter{相关工作}{Related Work}

\section{相关工作写作方式}{Writing Related Work}

相关工作章节应围绕论文问题组织，而不是简单罗列文献。可以按照研究主题、技术路线、应用场景或时间脉络划分小节，并在每个小节末说明已有工作与本文工作的关系。一个常见结构是“先介绍基础方法，再介绍与本文最接近的工作，最后指出本文的差异”。

\LaTeX{} 适合处理包含公式、交叉引用和参考文献的长文档\cite{lamport1994latex}。BibTeX 可以将参考文献信息与正文写作分离，便于维护统一的引用格式\cite{patashnik1988bibtexing}。

\section{方法类别一}{Method Category I}

这一节可以介绍最基础或最主流的一类方法。写作时应关注方法假设、核心思想和适用范围，而不是逐篇复述论文摘要。例如，可以按照“输入表示、模型结构、训练目标、评价方式”的顺序概括同一类方法的共同点。

当需要比较多篇工作时，可以使用如下句式：早期方法通常依赖某一假设，后续方法通过引入新的表示或训练目标缓解了该问题，但在某些场景下仍然存在限制。这样的写法比“某某提出了某方法，某某又提出了某方法”更容易形成论证。

\section{方法类别二}{Method Category II}

这一节可以介绍与本文更接近的一类方法。需要说明它们为什么相关，以及本文与它们的主要区别。区别可以来自任务设定、数据来源、模型结构、损失函数、实验协议或应用目标。

如果本文只是对已有方法进行复现、改进或应用，也应明确说明改动边界。例如：“本文不重新设计主干网络，而是在统一复现框架下比较不同训练策略。”这种表述可以降低过度声称的风险。

\section{评价协议与复现工作}{Evaluation and Reproducibility}

对于实验型论文，相关工作中也可以单独讨论评价协议。不同数据划分、预处理方式、随机种子和指标实现都可能影响结果。若后文实验依赖公开基线，应在此说明基线来源和可比性条件。

复现相关表述应尽量具体，例如说明是否使用官方代码、是否重新训练模型、是否只引用论文报告结果。若不同论文的实验设置并不一致，不宜直接把结果放入同一张表中进行强比较。

\section{与本文的关系}{Relation to This Thesis}

本节通常用于总结已有工作的不足，并说明本文的切入点。写作时应避免夸大已有方法的缺陷，也不要把“尚未有人做过”作为唯一动机；更稳妥的方式是指出具体任务、数据、方法假设或评价协议中的未解决问题。

一个可操作的收束方式是：已有工作已经解决了哪些问题，哪些问题仍未充分讨论，本文将在什么约束下尝试解决其中哪一部分。这样可以让读者在进入方法章节前清楚理解本文贡献的边界。
